% ----- CHAUPAI 5 -----

\newpage

\subsection*{\deva{पञ्चम चौपाई} (Fifth Chaupai)}
\addcontentsline{toc}{subsection}{\deva{पञ्चम चौपाई} (Fifth Chaupai) — \deva{हाथ बज्र...}}

\begin{minipage}[t]{0.55\textwidth}
\vspace{0pt}

{\large \linespread{1.5}\selectfont
\deva{हाथ बज्र औ ध्वजा बिराजै।}\\
\deva{काँधे मूँज जनेऊ साजै॥}
\par}

\vspace{0.5em}

{\large \linespread{1.5}\selectfont
\telu{హాథ బజ్ర ఔ ధ్వజా బిరాజై।}\\
\telu{కాఁధే మూఁజ జనేఊ సాజై॥}
\par}

\vspace{0.5em}

{\large \linespread{1.3}\selectfont
\textit{hātha bajra au dhvajā birājai |}\\
\textit{kāṃdhe mūṃja janeū sājai ||}
\par}
\end{minipage}%
\hfill
\begin{minipage}[t]{0.42\textwidth}
\vspace{0pt}
\begin{tcolorbox}[anchorbox]
\textbf{The Scholar and the Protector:} He holds a devastating weapon (\textit{Bajra}) in his hand, but wears the sacred thread (\textit{Janeu}) of a peaceful scholar on his shoulder. This is the ultimate balance we should all strive for: possessing a highly educated, peaceful mind, while remaining strong enough to fiercely protect others from injustice.
\vspace{0.5em}\newline He carried the absolute power of a warrior beneath the humble, unthreatening garb of a dedicated scholar.\newline\null\hfill\href{https://www.valmiki.iitk.ac.in/sloka?field_kanda_tid=4&language=dv&field_sarga_value=3}{\textit{Kishkindha Kanda, 3:2}}
\end{tcolorbox}
\end{minipage}

\vspace{0.5em}
\subsubsection*{\textbf{Visual Rhythm Map:}}
\begin{table}[h!]
\centering
\renewcommand{\arraystretch}{1.2}
\setlength{\tabcolsep}{0pt}
\begin{tabularx}{\textwidth}{|*{16}{>{\centering\arraybackslash\hsize=1.0\hsize}X|}}
\hline
\multicolumn{3}{|>{\centering\arraybackslash\hsize=3.0\hsize}X|}{\textbf{\guru{हा}\laghu{थ}} \par (3)} &
\multicolumn{3}{>{\centering\arraybackslash\hsize=3.0\hsize}X|}{\textbf{\guru{ब}\laghu{ज्र}} \par (3)} &
\multicolumn{2}{>{\centering\arraybackslash\hsize=2.0\hsize}X|}{\textbf{\guru{औ}} \par (2)} &
\multicolumn{3}{>{\centering\arraybackslash\hsize=3.0\hsize}X|}{\textbf{\laghu{ध्व}\guru{जा}} \par (3)} &
\multicolumn{5}{>{\centering\arraybackslash\hsize=5.0\hsize}X|}{\textbf{\laghu{बि}\guru{रा}\guru{जै}} \par (5)} \\
\hline
\end{tabularx}
\vspace{-1.5em}
\end{table}

\begin{table}[h!]
\centering
\renewcommand{\arraystretch}{1.2}
\setlength{\tabcolsep}{0pt}
\begin{tabularx}{\textwidth}{|*{16}{>{\centering\arraybackslash\hsize=1.0\hsize}X|}}
\hline
\multicolumn{4}{|>{\centering\arraybackslash\hsize=4.0\hsize}X|}{\textbf{\guru{काँ}\guru{धे}} \par (4)} &
\multicolumn{3}{>{\centering\arraybackslash\hsize=3.0\hsize}X|}{\textbf{\guru{मूँ}\laghu{ज}} \par (3)} &
\multicolumn{5}{>{\centering\arraybackslash\hsize=5.0\hsize}X|}{\textbf{\laghu{ज}\guru{ने}\guru{ऊ}} \par (5)} &
\multicolumn{4}{>{\centering\arraybackslash\hsize=4.0\hsize}X|}{\textbf{\guru{सा}\guru{जै}} \par (4)} \\
\hline
\end{tabularx}
\vspace{-1.5em}
\end{table}

\vspace{1em}
\textbf{ \deva{सरल-संस्कृतम्}:}\\
 \deva{भवतः हस्ते वज्रं ध्वजः च विराजते।}\\
 \deva{भवतः स्कन्धे मौञ्जी-यज्ञोपवीतं शोभते॥}\\

In your hand shines the thunderbolt and the banner.\\
On your shoulder, the sacred thread of Munja grass looks beautiful.

\vspace{1em}
\newpage
\textbf{\deva{प्रतिपदार्थः} (Meaning of each word):}
\begin{table}[h!]
\centering
\renewcommand{\arraystretch}{1.2}
\begin{tabularx}{\textwidth}{@{} l l X X @{}}
\toprule
\textbf{Awadhi} & \textbf{Sanskrit} & \textbf{English} & \textbf{Grammar \& Notes} \\
\midrule
\deva{हाथ} / \telu{హాథ}  &  \deva{हस्ते}  &  In hands  &   \\
\deva{बज्र} \gotchaicon / \telu{బజ్ర}  &  \deva{वज्रम्}  &  Thunderbolt  & \\ 
\deva{औ} / \telu{ఔ}  &  \deva{और / च}  &  And  &   \\
\deva{ध्वजा} / \telu{ధ్వజా}  &  \deva{ध्वजः}  &  Flag / Banner  &   \\
\deva{बिराजै} / \telu{బిరాజై}  &  \deva{विराजते}  &  Shines  &   \\
\deva{काँधे} \gotchaicon / \telu{కాఁధే} & \deva{स्कन्धे} & On the shoulder & \\
\deva{मूँज} \gotchaicon / \telu{మూఁజ} & \deva{मौञ्जी} & Munja grass & Moon-dot (Chandrabindu) \\
\deva{जनेऊ} / \telu{జనేఊ}  &  \deva{यज्ञोपवीतम्}  &  Sacred thread  &   \\
\deva{साजै} / \telu{సాజై}  &  \deva{सज्जते / शोभते}  &  Adorns / Looks beautiful  &   \\
\bottomrule
\end{tabularx}
\end{table}

\begin{tcolorbox}[gotchabox]
\begin{itemize}
    \item \textbf{Nasalization:} Pay close attention to the \deva{चन्द्रबिन्दु} (Chandrabindu / Moon-dot) on \deva{काँधे} (Kaandhe) and \deva{मूँज} (Moonja). This nasalizes the vowel without adding the extra weight of a full consonant, preserving the 16-beat meter.
    \item \textbf{Conjunct Consonant Weight (\deva{बज्र}):} The `a` in `ba` becomes a 2-beat Guru because it precedes the conjunct `jra`!
\end{itemize}
\end{tcolorbox}

\newpage
